\chapter{模板计算与线程粗化}

\begin{solutionexercise}{1}
\problemheading 考虑在包含边界单元的 $120\times120\times120$ 网格上执行三维模板计算。
\begin{enumerate}[label=\alph*.,leftmargin=2.2em]
  \item 每次模板扫掠会计算多少个输出网格点？
  \item 对图 8.6 的基本内核，假设线程块大小为 $8\times8\times8$，
  需要多少个线程块？
  \item 对图 8.8 的共享内存分块内核，假设线程块大小为
  $8\times8\times8$，需要多少个线程块？
  \item 对图 8.10 同时采用共享内存分块和线程粗化的内核，假设线程块
  大小为 $32\times32$，需要多少个线程块？
\end{enumerate}
沿用本章一阶三维七点模板：只更新每维下标 $1\ldots N-2$ 的内部点。图 8.6 的
$8^3$ 块直接覆盖全网格；图 8.8 的 $8^3$ 输入瓦片产生 $6^3$ 输出瓦片；图 8.10
的 $32\times32$ 输入平面产生 $30\times30$ 输出平面，并在 $z$ 方向连续粗化
处理 30 个输出平面。

\approachheading 一阶七点模板只更新每维去掉两侧各一层后的内部点。基本内核按 8
覆盖全网格；共享内存内核的 $8^3$ 输入块产生 $6^3$ 输出；$32^2$ 粗化块每维有效
输出边长为 30，并沿 $z$ 连续处理 30 个平面。

\answerheading
\begin{enumerate}[label=\alph*.,leftmargin=2.2em]
\item 输出点数为
\[
(120-2)^3=118^3=1{,}643{,}032\ \text{个网格点}.
\]
\item 基本内核每维需 $\lceil120/8\rceil=15$ 个块，共
$15^3=3375$ 个块。
\item 原书图 8.8 中输入瓦片边长 8、输出瓦片边长 $8-2=6$，每维需
$\lceil118/6\rceil=20$ 个块，共 $20^3=8000$ 个块。
\item 原书图 8.10 中 $32\times32$ 输入平面对应 $30\times30$ 输出平面，并沿
$z$ 粗化 30 层。每维需 $\lceil118/30\rceil=4$ 个块，共 $4^3=64$ 个块。
\end{enumerate}
边界块中的越界或边界线程由内核条件屏蔽，不改变启动块数。

\complexity{$O(118^3)$ 有效工作}{基本内核 $O(1)$；分块内核使用片上瓦片}{输出点在三维块网格中并行}
\conclusionheading (a) 1,643,032；(b) 3375；(c) 8000；(d) 64。
\discussionheading
\discussionpoint{七点模板把离散物理写成局部数据依赖。}
中心点与六个轴向邻点的加权和看似一种小卷积，但这些系数通常来自偏微分方程的有限差分，例如三维扩散或 Poisson 方程，而不是任意图像滤波器。固定边界单元不更新，等价于把边界值作为 Dirichlet 条件；若改成周期边界，越过左面会从右面取值，Neumann 条件则可能通过镜像或 ghost 值表达导数。边界规则改变后，有效输出点数、邻域索引和数值解的物理含义都会改变。模板优化因此不能只验证 GPU 与某段代码一致，还应检查离散方程的守恒、收敛或边界不变量。
\discussionpoint{启动线程、加载线程和输出线程是三种角色。}
基本内核可为整个 $120^3$ 体积启动线程，再让边界线程跳过更新；分块内核还可能让外圈线程只负责把 halo 载入共享内存，而不生成输出。于是“启动了多少线程”不能直接换算成“计算了多少网格点”，更不能单独判断效率。一个 tile 的额外加载如果让许多相邻输出复用同一全局值，可能比每线程都计算却重复读取更划算；反之，小 tile 的 halo 面积占比很高，协作开销可能吞掉收益。应分别统计有效输出、边界屏蔽、协作加载和相邻块重复 halo，才能把线程冗余与数据复用放在同一账本中比较。
\discussionpoint{时间分块把复用从空间推进到迭代维。}
共享内存空间分块只复用当前时间步的邻点，而 Jacobi 等迭代会连续读取下一时刻刚生成的网格；若在片上推进多个时间步，就能避免每一步都把整个体积写回并重新读取。代价是依赖锥随时间扩张：计算 $k$ 步后的内部点需要初始 tile 周围厚度为 $k$ 的 halo，每推进一步可安全输出的内部区域都会收缩。还要保存多个时间层或采用交替缓冲，并确保同步不会混用新旧状态。时间分块因此提升算术强度的同时加剧容量、寄存器和边界复杂度，它是依赖图变换而不只是多循环几次。
\discussionpoint{多 GPU 模板受表面积与体积比例支配。}
把三维域切成子域后，每个 GPU 主要计算自己的体积，却在每轮交换六个面上的 halo；计算量随边长立方增长，通信量近似随边长平方增长，所以较大的子域更容易摊薄通信。子域过小或切分方向不合适时，网络延迟和边界复制会压过单 GPU 内核优化，即使局部 tile 的带宽表现很好。工程实现常把内部点计算与 halo 传输重叠，再在数据到达后处理边界层，但这要求明确 stream、通信库和设备拓扑依赖。评价结果应同时报告每秒网格更新、通信隐藏比例与数值收敛一致性，才能连接局部内核吞吐和整体求解时间。
\variantheading 若网格为 $62^3$，内部点为 $60^3=216000$；$8^3$ 基本块需 $8^3=512$ 块。
\practiceheading 用 Nsight Compute 比较三内核的 DRAM Bytes、active threads 和 achieved occupancy。
\sourcecheck{https://docs.nvidia.com/cuda/cuda-c-programming-guide/}{NVIDIA CUDA C++ Programming Guide 的三维网格向上取整规则；并按原书图 8.6、8.8、8.10 的索引独立复核}
\end{solutionexercise}

\begin{solutionexercise}{2}
\problemheading 考虑一个同时采用共享内存分块和线程粗化的三维七点模板实现。该实现与图
8.10 和图 8.12 类似，但分块不是规则立方体；它采用
$32\times32$ 的线程块，并把粗化因子设为 16，也就是说，每个线程块沿 $z$ 维
处理连续的 16 个输出平面。
\begin{enumerate}[label=\alph*.,leftmargin=2.2em]
  \item 在线程块的整个生命周期中，它所加载的输入分块大小是多少（以元素个数
  计）？
  \item 在线程块的整个生命周期中，它所处理的输出分块大小是多少（以元素个数
  计）？
  \item 该内核的算术强度是多少（以 FLOP/B 计）？
  \item 如果不使用寄存器分块，如图 8.10 所示，每个线程块需要多少
  字节共享内存？
  \item 如果使用寄存器分块，如图 8.12 所示，每个线程块需要多少
  字节共享内存？
\end{enumerate}
每个网格元素为 4 B 单精度浮点数；七点模板每个输出执行 7 次乘法和 6 次加法。
输入在 $x,y$ 两维各含一层 halo，在 $z$ 向的 16 个输出平面前后也各含一个输入平面。
算术强度按每个输入元素从全局内存加载一次、每个输出元素写回一次计，忽略系数流量
与边界块浪费。图 8.10 同时在共享内存保存 prev/curr/next 三个平面，图 8.12 只在
共享内存保存 curr 平面而把另外两个平面保存在寄存器。

\approachheading 一阶模板在 $x,y$ 两维各保留一层 halo，所以每个输入平面为
$32\times32$、输出平面为 $30\times30$；16 个输出平面在 $z$ 向还需前后各一个
输入平面。

\answerheading
\begin{enumerate}[label=\alph*.,leftmargin=2.2em]
\item 生命周期中加载的输入瓦片为 $32\times32\times(16+2)$，共
\[
32^2\times18=18{,}432\ \text{个元素}.
\]
\item 输出瓦片为 $30\times30\times16$，共
$14{,}400$ 个元素。
\item 每个输出执行 13 FLOP（7 乘、6 加）。全局流量按每个输入加载一次、每个输出
存储一次计，因此
\[
I=\frac{13\times14{,}400}
        {4(18{,}432+14{,}400)}
 =\frac{325}{228}\approx1.42544\ \mathrm{FLOP/B}.
\]
\item 不使用寄存器分块时同时保存 prev/curr/next 三个输入平面，共
$3\times32^2\times4=12{,}288$ B，即 12 KiB 共享内存。
\item 使用寄存器分块后，prev/next 留在线程寄存器，只需共享当前平面，共
$32^2\times4=4096$ B，即 4 KiB。
\end{enumerate}
算术强度忽略系数流量、缓存重复和边界块浪费，正是题目采用的理想化模型。

\complexity{$13\times14{,}400$ FLOP/块}{12 KiB 或 4 KiB 共享内存/块}{1024 个线程协作处理 16 个输出平面}
\conclusionheading (a) 18,432；(b) 14,400；(c) 约 1.42544 FLOP/B；
(d) 12,288 B；(e) 4096 B。
\discussionheading
\discussionpoint{沿 z 粗化形成一个寄存器滑动窗口。}
若同一线程连续计算多个 z 位置，当前输出所需的“前、当前、后”三个中心值会在下一输出中平移角色：旧当前成为新前，旧后成为新当前，只需再载入一个新后值。把这三个值留在寄存器中，就不必为每个 z 输出重新经过共享内存，类似 CPU 模板代码中的滑动窗口或流式管线。复用成立的前提是线程始终沿同一 x、y 坐标前进，并严格区分尚未进入、正在使用和已经离开的平面。这个变换把数据局部性从块内线程协作提升为线程私有时间序列，也解释了粗化为何能减少某些片上访问。z 边界的 ghost 值还必须在窗口推进前定义，不能让首尾迭代读取未初始化寄存器。
\discussionpoint{不同方向的邻点适合不同通信范围。}
x、y 邻点属于同一输出平面，一个线程始终可以按自己的邻点下标直接从全局内存重载它们；缓存和合并访问可能使这条朴素路径已经足够有效。只有当设计意图是复用相邻线程已经装入寄存器的值、避免再次访问全局内存时，才需要共享内存，或在活动掩码与线程映射允许时使用 warp shuffle 等线程交换机制。z 邻点来自同一线程即将连续处理的位置，因此更适合保存在寄存器滑动窗口中。优化方案可把当前 x/y 平面放进共享内存、把前后中心列留在线程私有状态，但这是一种复用选择而非正确性的强制条件；共享容量下降还会与滚动变量延长寄存器活跃区间相互制约。应依据重载流量、线程作用域和实际缓存行为选择存储层次。
\discussionpoint{理想算术强度会被粗化尾部和物理流量修正。}
增加 z 粗化因子可以让上下 halo 平面由更多内部输出摊薄，所以按“每个输入只从 DRAM 载入一次、每个输出写一次”的模型，单位输出字节数逐渐下降。然而线程数量也随每线程输出数增加而减少，最后一个 z 分段若不足粗化长度，会出现闲置循环与边界谓词。缓存行过取、未合并的跨平面地址、寄存器 spill 以及相邻块重复 halo 都会让实测 DRAM 字节偏离公式。Roofline 上的理论点只是上界假设，只有把实际传输量代回算术强度，才能判断优化真正移动了数据屋顶还是仅改变了源码计数。
\discussionpoint{调优应搜索形状组合并守住数值更新语义。}
x/y tile、z 粗化长度、共享平面数量和是否双缓冲会共同决定线程数、共享容量、寄存器及同步次数，单独扫描一个参数容易错过资源间的阶梯变化。可以先用编译器资源报告排除发生 spill 的组合，再在目标尺寸上比较实际流量和时间，并用 CPU 参考覆盖非整除边界。若进一步跨多个迭代做时间分块，依赖锥会扩大，内部有效区域收缩，而且 Jacobi 的新旧数组语义不能无意中变成 Gauss--Seidel 的就地更新。性能搜索必须与守恒量、残差或收敛曲线一起验证，才能确保更快的内核仍在求解同一个离散问题。
\variantheading 粗化因子改为 8 时，输入为 10240、输出为 7200 个元素，强度约 1.3417 FLOP/B。
\practiceheading 用 ptxas 和 Nsight Compute 同时检查寄存器数、shared bytes、spill 与 Roofline 位置。
\sourcecheck{https://docs.nvidia.com/cuda/cuda-c-best-practices-guide/}{NVIDIA CUDA Best Practices Guide 的共享内存与算术/带宽计量原则；元素数、FLOP 和字节三路独立复核}
\end{solutionexercise}
